\scanpage{141}
\mspage{- 129 -}
% archive sheet number: 140
\hfill n\textsuperscript{o} 316

où, dans la série $o_3$, tous les monômes sont de degré $\handwritten{\geqslant}3$, par rapport à l'ensemble des $Y_i$ et des $Z_i$. En déduire que :
\[
\begin{aligned}
u(f(x,y))-u(f(y,x))&=\handwritten{\sum}_{i\handwritten{<}j}(X_iY_j-X_jY_i)([T_i,T_j]^{(e)}(u))+\\
&\qquad +o'_3(x,y)
\end{aligned}\tag*{(3)}
\]
où tous les termes de $o'_3$ sont de degré total $\handwritten{\geqslant}3$. Montrer que si $G$ est commutatif, $g$ est commutative.

d) Soit $G'$ un second groupe formel, de dimension $n$, dont la loi de groupe est $f'$. On appelle \uline{homomorphisme formel} de $G$ dans $G'$ un système $F=(F_j(x))_{1\handwritten{\leqslant}j\handwritten{\leqslant}m}$ d'éléments de $A_x$ tel que l'on ait $f'(F(x),F(y))=F(f(x,y))$. Pour toute dérivation invariante à gauche $D\handwritten{\in}g$, montrer que $v\handwritten{\longrightarrow}D(v\unsure{o|\circ}F)$ est une dérivation invariante à gauche appartenant à l'algèbre de Lie $g'$ de $G'$. Si on la désigne par $F^{\handwritten{*}}(D)$, $F^*$ est un homomorphisme (un $p$-homomorphisme pour $p>0$) de $g$ dans $g'$. Si $(T_i)_{1\handwritten{\leqslant}i\handwritten{\leqslant}n}$, $(T'_j)_{1\handwritten{\leqslant}j\handwritten{\leqslant}m}$ sont les bases de $g$ et $g'$ telles que $T_i^{(e)}=D_i^{(e)}$ et $T'_j{}^{(e)}=D_j^{(e)}$, montrer que la matrice de $F^{\handwritten{*}}$, rapportée à ces bases, est la matrice $(\handwritten{\partial}F_j/\handwritten{\partial}X_i)_c$ (terme constant de $\handwritten{\partial}F_j/\handwritten{\partial}X_i$, où $i$ est l'indice des lignes, $j$ celui des colonnes).

25) - On appelle \uline{groupe linéaire formel} à $n$ variables sur le corps $K$, et on désigne par $\mathrm{GL}(n,K)$ (lorsqu'aucune confusion n'est possible) le groupe formel de dimension $n^2$ dont la loi de groupe est définie par $f_{ij}(u,v)=-\handwritten{\delta}_{ij}+\handwritten{\sum}_{k=1}^{n}(\handwritten{\delta}_{ik}+U_{ik})(\handwritten{\delta}_{kj}+V_{kj})$ ($\handwritten{\delta}_{ij}$ indice de Kronecker ; les $3n^2$ indéterminées de la définition générale de l'exerc. 24 sont ici notées $U_{ij}$, $V_{ij}$, $W_{ij}$ avec 2 indices variant de 1 à $n$). Montrer que, si $D_{ij}=\handwritten{\partial}/\handwritten{\partial}U_{ij}$, les dérivations
